\documentclass[11pt]{article}

\usepackage[a4paper,margin=1in]{geometry}
\usepackage[T1]{fontenc}
\usepackage[utf8]{inputenc}
\usepackage{lmodern}
\usepackage{enumitem}
\usepackage{hyperref}

\begin{document}

\begin{flushleft}
11-Apr-2026\\[0.75em]
Manuscript ID: BIOINF-2026-0316\\
Title: \textit{PyFgsea: A Rust-Powered, fgseaMultilevel-Aligned GSEA Framework with Rolling-Window Enrichment Along Single-Cell Trajectories}
\end{flushleft}

\vspace{0.5em}

Dear Associate Editor Russell Schwartz and Reviewers,

We sincerely thank you for the careful evaluation of our manuscript and for the opportunity to submit a substantially revised version. The reviewers' comments made clear that our original submission did not explain the implementation boundaries, validation scope, and practical limitations with sufficient clarity in the main text. We therefore substantially rewrote both the manuscript and the Supplementary Material with four priorities in mind: improving main-text self-contained clarity, clarifying methodological and implementation boundaries, strengthening rolling-window robustness and interpretation, and ensuring archival and format compliance.

\section*{Executive Summary for the Associate Editor}

\begin{enumerate}[leftmargin=1.5em]
\item \textbf{Main-text self-contained clarity strengthened.} We rewrote Sections~2.1--2.3 and 3.1--3.3 so that the main manuscript now explicitly states the relationship to \texttt{fgseaMultilevel}, includes compact ES and multilevel tail-factorization formulas, reports a concrete large-scale benchmark example (runtime and PeakRSS), and summarizes ES/NES/transformed nominal-$P$ agreement more directly. We also now state explicitly that nominal $P$-value agreement can be weaker in sparse local trajectory windows than ES or NES agreement.

\item \textbf{Methodological transparency and implementation boundaries clarified.} The revised manuscript now states directly that PyFgsea is not a new ES statistic or a new rare-event estimator. Rather, it re-implements the \texttt{fgseaMultilevel} rare-event estimator target while differing at the implementation layer (Rust/PyO3 backend, Rayon scheduling, deterministic pathway-local PRNG streams) and extending the workflow with a stateful rolling-window runner for trajectory analysis. The Supplementary Material now adds the requested workflow diagram, explicit formulas, a same-versus-different table, and a pseudocode-style algorithm box.

\item \textbf{Rolling-window robustness and interpretation strengthened.} We expanded the trajectory-related material by adding empirical parameter guidance in the main text, a fuller sensitivity analysis across multiple pathways and window/step combinations, a boundary-handling audit, and an explicit explanation of why the illustrative erythroid example uses a 500-cell baseline window (approximately 14\% of that trajectory) even though our generic starting recommendation is 5--10\%. We also unified several statements that previously risked ambiguity: within-window significance now explicitly means Benjamini--Hochberg correction across pathways, Figure~2 is stated to use the default full-window workflow, and rolling-window acceleration is reported as approximately 1.9-fold end-to-end in a conservative stress test and up to 7.47-fold in a narrower component benchmark.

\item \textbf{Archival and format compliance completed.} To comply with the Application Note format, the revised main text now contains exactly two illustrative items (Figure~1 and Figure~2), with the original performance table moved to Supplementary Table~S4. We also archived a permanent snapshot of the software on Zenodo and added the DOI to the Availability statement: \href{https://doi.org/10.5281/zenodo.19446446}{10.5281/zenodo.19446446}.
\end{enumerate}

\section*{Clarifications on Changes Not Made}

\begin{enumerate}[leftmargin=1.5em]
\item \textbf{We did not replace the core statistical estimator.} We did not introduce a new ES definition or a different rare-event estimator, because the purpose of PyFgsea is to provide a Python/Rust implementation aligned with the \texttt{fgseaMultilevel} rare-event estimator target rather than propose a new enrichment statistic. We instead made this boundary much more explicit in the revision.

\item \textbf{We did not move every validation grid into the main text.} We did, however, move compact formulas, concrete benchmark numbers, and stronger quantitative summaries into the main manuscript. The Supplementary Material remains the appropriate location for the full validation grids, thread audits, and parameter stress tests needed to document implementation robustness without exceeding the journal's format constraints.
\end{enumerate}

Below we provide a point-by-point response to the reviewers.

\section*{Reviewer 1}

\begin{enumerate}[leftmargin=1.5em]
\item \textbf{Justification and sensitivity analysis for rolling-window parameters; suggest empirical starting points.}

\textbf{Response:} We thank the reviewer for this practical and important suggestion. We agree with the reviewer that the original parameter justification was too limited. In the revised main text (Section~2.3), we now present approximately 5--10\% of total cells for the window size and 1--2\% of total cells for the step size as a practical starting range that should be tuned to trajectory length and noise level. We also clarified that the erythroid illustration in Figure~2 intentionally uses a somewhat larger 500-cell baseline window (approximately 14\% of that analyzed trajectory) because this lineage segment is relatively short and noisy, and the figure was intended to emphasize a smoother descriptive profile rather than maximal temporal resolution. To avoid overstating the evidence, the revised Supplementary text now also notes that the current erythroid sensitivity grid brackets this generic starting range rather than directly centering on it.

In the Supplementary Material, we substantially expanded the sensitivity analysis. Figure~S8 now shows both window-size and step-size sensitivity across three representative pathways. Figure~S13 and Table~S9 summarize the full window-by-step grid using smoothness, peak-shift, significant-window overlap, and runtime diagnostics. Figure~S14 further audits terminal boundary handling explicitly. We believe these additions make the parameter trade-offs much more transparent to users.

\item \textbf{Discussion on limitations of the rolling-window ranking statistic.}

\textbf{Response:} We fully agree that these limitations should be stated explicitly. We therefore added a dedicated paragraph to the Discussion section explaining that the rolling-window difference-of-means ranking statistic assumes locally monotonic expression shifts and can therefore underestimate complex non-monotonic dynamics, leading to false negatives. We also note the converse risk that sparse genes with high technical variance may dominate local rankings and create false positives. This addition is intended to support more cautious and appropriate interpretation of rolling-window results.
\end{enumerate}

\section*{Reviewer 2}

\begin{enumerate}[leftmargin=1.5em]
\item \textbf{Lack of methodological workflow diagram.}

\textbf{Response:} We appreciate this point. We added a dedicated methodological workflow diagram as Supplementary Figure~S7, and we now refer to it explicitly in Section~2.1 of the main text.

\item \textbf{Insufficient description of core workflows and missing core formulas.}

\textbf{Response:} We agree that the original version omitted too much mathematical detail. We addressed this at two levels. First, the revised main text now includes a compact ES formula and a compact multilevel tail-factorization formula in Section~2.1, so that the central estimator is no longer described only by reference to the Supplementary Material. Second, Supplementary Section~S8 was rewritten to give a fuller algorithmic description, including the explicit definitions of $P_{\mathrm{hit}}(i)$ and $P_{\mathrm{miss}}(i)$, the random pathway generation strategy, and a pseudocode-style summary of the multilevel routine as implemented in the current Rust release.

\item \textbf{Lack of detail in multilevel adaptive sampling, including specific differences from the original method.}

\textbf{Response:} We thank the reviewer for emphasizing this point. In the revised manuscript, the main text now states directly that PyFgsea re-implements the \texttt{fgseaMultilevel} rare-event estimator target with the same ES definition, null-pathway construction, and multilevel tail factorization, while also clarifying that the current Rust release realizes the conditional levels through empirical median thresholding and elite-clone rejuvenation. We then distinguish implementation-level differences from methodological continuity. In Supplementary Table~S5, we explicitly separate: (i) components preserved from \texttt{fgseaMultilevel}, (ii) implementation-level differences in PyFgsea (compiled Rust backend, work-stealing scheduling, deterministic pathway-local PRNG streams), and (iii) the PyFgsea-specific rolling-window extension. Supplementary Section~S8 now gives a more explicit pseudocode-style description of initialization, empirical level thresholds, elite selection, constrained rejuvenation steps, and stopping rules.

\item \textbf{Inadequate description of multi-threading and unspecified parallel random number management.}

\textbf{Response:} We expanded this substantially. Section~2.2 of the main text now states that PyFgsea parallelizes pathway-level evaluations using Rayon dynamic work stealing and that it guarantees deterministic reproducibility under matched inputs and a fixed master seed by deriving independent pathway-local PRNG streams from the user-provided master seed and the unique pathway identifier. Supplementary Section~S9 provides the longer architectural explanation. To make this claim empirical rather than purely descriptive, we also added a fixed-seed cross-thread audit in Supplementary Section~S12 (Figure~S11 and Table~S8), showing that pathway-wise outputs and pathway ranking remained invariant across 1--16 threads under matched inputs and a fixed master seed.

\item \textbf{Lack of quantitative metrics for equivalence validation.}

\textbf{Response:} We agree that our original qualitative language was insufficient. Section~3.2 of the revised main text now reports explicit quantitative agreement metrics for NES (Pearson $r=1.000$, Spearman $\rho=1.000$, RMSE $=0.0135$, median absolute difference $|\Delta|=0.0070$) and also summarizes transformed nominal-$P$ concordance. We also revised the wording to be more candid: ES was numerically identical within machine precision, NES remained near-identical, whereas transformed nominal $P$-values were overall concordant but could show visibly larger dispersion in sparse local trajectory windows. The expanded cross-regime summaries are now provided in Supplementary Figure~S10 and Tables~S6--S7.

\item \textbf{Overall concern that the manuscript sacrificed too many technical details and the main text was not sufficiently self-contained.}

\textbf{Response:} We understood this as the central criticism, and we tried to address it directly rather than only by adding more supplementary experiments. The revised main text now states the estimator relationship to \texttt{fgseaMultilevel} explicitly, includes compact ES and multilevel tail-factorization formulas, reports a concrete large-benchmark runtime/memory example, summarizes ES/NES/transformed nominal-$P$ agreement quantitatively, clarifies that within-window FDR is controlled across pathways, and adds more explicit parameter guidance plus a limitations paragraph for rolling-window analysis. The Supplementary Material now carries the workflow diagram, the full algorithm box, the detailed same-versus-different table, and the larger validation grids. This revised division of labor is intended to make the main text materially more self-contained while still respecting the Application Note format.
\end{enumerate}

We believe these revisions substantially address the reviewers' concerns while respecting the Application Note format constraints.

\vspace{1em}

Sincerely,\\
Kuanghao Wang and Hong Shi

\end{document}
